\documentclass[10pt,a4paper]{amsart}
\usepackage[margin=19mm]{geometry}
\usepackage{amsmath,amssymb,mathtools}
\usepackage[scr=rsfs,cal=euler]{mathalfa}
\usepackage{xfrac}
\usepackage[unicode,hidelinks,bookmarks=false]{hyperref}
\newcommand{\ie}{i.e.\ }
\newcommand{\cf}{cf.\ }
\newcommand{\resp}{respectively\ }
\newcommand{\Subsection}{\subsection}
\providecommand{\nobreakdash}{\nobreak\hbox{-}}
\setlength{\emergencystretch}{2em}
\multlinegap0pt
\usepackage{bm}
\usepackage{bbm}
\usepackage{dsfont}
\usepackage{xspace}
\usepackage{cancel}
\usepackage{mathtools}
\usepackage{amsthm}
\makeatletter
\@ifundefined{thm}{\newtheorem{thm}{Thm}}{}
\@ifundefined{prop}{\newtheorem{prop}[thm]{Proposition}}{}
\makeatother
\newcommand{\N}{\mathbb{N}}
\newcommand{\W}{\mathcal{W}}
\title[Centralizer Rigidity near Elements of the Weyl Chamber Flow]{Centralizer Rigidity near Elements of the Weyl Chamber Flow}
\author{Zhijing Wendy Wang}
\date{Source: arXiv:2309.07282v3 (CC BY 4.0)}
\begin{document}
\maketitle

\noindent\emph{Source and licence.} Centralizer Rigidity near Elements of the Weyl Chamber Flow. Version arXiv:2309.07282v3 (CC BY 4.0), distributed under the Creative Commons Attribution 4.0 International licence, as recorded in the arXiv licence metadata for this exact version and shown on the abstract page https://arxiv.org/abs/2309.07282v3. Full rights notice: licenses/arxiv-2309.07282-CC-BY-4.0.txt.\par\smallskip

\noindent\textit{Editorial scope.} Counted unit: the Prop whose source label is \texttt{None} in the article (source0.tex lines 494--496, source label \texttt{None}), delivered together with the article's own complete original proof of it (source0.tex lines 498--519, one \texttt{proof} environment of 22 lines). No mathematics is written, repaired or re-derived: statement and proof are the article's own text line for line. Required context retained uncounted: none. Every definition, display and label of those lines is retained unchanged. Omitted from this package and not redistributed: 1--493, 497--497, 520--1634. Nothing inside a retained excerpt is omitted or altered: every retained body is delivered exactly as the article's lines are written, so no omission receipt of any kind accompanies this package. Citations: the retained excerpts carry no citation command at all, so no bibliography entry of the article is transcribed and no reference is redistributed. Reference adaptation: none was needed and none was made; every reference inside the delivered unit resolves inside it, so no unresolved reference or citation is produced. Printed slips kept literal: no printed word, symbol, hypothesis or number of the counted statement or of its proof was altered. Layout and class: the article's own class is \texttt{amsart} with packages including inputenc, amsmath, amssymb, amsfonts, amsthm, graphicx, hyperref, enumerate, tikz; the wrapper is the lane's pinned \texttt{amsart} editorial wrapper at 10pt on A4, which restates the theorem-like environments the retained text needs and adds the article's own macro definitions that the retained text needs (\texttt{\textbackslash N}, \texttt{\textbackslash W}), with extra wrapper packages (bm, bbm, dsfont, xspace, cancel, mathtools). The delivered numbering is the unit's own. Assets: no retained span contains a figure environment, a graphics-inclusion command, a PDF-page inclusion, a table or a TikZ picture; no image file is copied into this package, the package declares no asset dependency and no image file is redistributed with it. Licence: version arXiv:2309.07282v3 of this article is distributed under the Creative Commons Attribution 4.0 International licence, as recorded in the arXiv licence metadata for this exact version and shown on the abstract page https://arxiv.org/abs/2309.07282v3. A full rights notice is delivered at licenses/arxiv-2309.07282-CC-BY-4.0.txt. Characters: every retained excerpt is pure ASCII, so the delivered engine, XeLaTeX with its default Latin Modern text font, reports no missing character.
\begin{prop}
 Let $f_0$ be a non-trivial element of the Weyl chamber flow, then any $C^1$-perturbation $f$ of $f_0$ has global holonomy in the universal cover. In particular, $f$ has global holonomy in a center leaf $\W^c_f(x_0)$ if $\W^c_f(x_0)$ is simply connected.
\end{prop}

\begin{proof}
 We lift the diffeomorphisms $f,f_0$ and the corresponding foliations in $G/\Gamma$ to the cover $ G$, denoted by $\tilde{f_0},\tilde f:G\to G$ and $\tilde{\mathcal W}^*_{f_0}$, $\tilde{\mathcal W}^*_f$, for $*=s,c,u,sc,cu$.

 Since $f_0$ is an affine diffeomorphism, $\tilde{\mathcal W}^*_{f_0}$ are exactly the orbits of $G^*=\exp(\mathfrak{g}^*)$ under left translation, where $\mathfrak{g}=\mathfrak{g}^s\oplus \mathfrak{g}^c\oplus \mathfrak{g}^u$ is the invariant splitting into generalized eigenspaces with eigenvalues $|\lambda|<1$, $|\lambda|=1$ and $|\lambda|>1$, given by the action of $Df_0$ on the Lie algebra $\mathfrak{g}$; and $\mathfrak{g}^{cs}=\mathfrak{g}^{c}\oplus \mathfrak{g}^{s},\mathfrak{g}^{cu}=\mathfrak{g}^{c}\oplus \mathfrak{g}^{u}$.
 
 We first claim that the foliations $\tilde\W^c_{f_0}$ and $\tilde\W^s_{f_0}$ have global product structure inside a leaf $\tilde \W^{cs}_{f_0}$. This is because we have a natural map $G^c\times G^s\to G^{cs},(g_c,g_s)\mapsto g_cg_s$ which is a diffeomorphism.

 This allows us to define a continuous map on each $cs$-leaf of $\tilde f_0$ given by $\tilde \W^{cs}_{f_0}(x)\to G^s: y\mapsto g_s(yx^{-1})$ where $g_s: G^{cs}\to G^s$ is just the projection under the product structure $g_cg_s\mapsto g_s$.

 Since $f_0$ is dynamically coherent, normally hyperbolic, center bunched with smooth center leaves, for sufficiently small $C^1$-perturbations $f$ of $f_0$, there exists a leaf conjugacy $\phi$ from $(f_0,\W^c_{f_0})$ to $(f,\W^c_{f})$. By the dynamics of $f$ and $f_0$, $\phi$ also sends $\W^{cs}_{f_0}$ to $\W^{cs}_{f}$. We take the map $u:\tilde\W^{cs}_f(x_0)\to G^s$ given by $u(x)=g_s(\phi^{-1}(x)\cdot (\phi^{-1}(x_0))^{-1})$. Then by construction, leaves of $\tilde\W^c_f$ are exactly the level sets of the map $u$.

 We claim that $u$ restricted to any stable leaf $L=\tilde\W^s_f(x_1)$ is a homeomorphism. 
 
 Firstly, we show that $u|_L$ is injective. If $u(x)=u(y)$ for $x,y\in \tilde \W^s_f(x_1)$, then $y\in \tilde \W^c_f(x)\cap \tilde \W^s_f(x)$. By the dynamics in the stable and center foliation, this shows that $d(f^n(x),f^n(y))\le C_1 \mu^n d(x,y)$ and $d_{\tilde \W^c_f}(f^n(x),f^n(y))\ge C_2 \nu^n d_{\tilde \W^c_f}(x,y)$ for any $n>0$, where $0<\mu<\nu<1, C_1,C_2$ are fixed constants given by the dominated splitting that depends only on $f_0$ and the $C^1$-neighborhood where we pick $f$. Since the leaf conjugacy $\phi^{-1}$ is $C^1$-close to identity along $\W^c_f$ and $C^0$-close to identity in $X$, we have $$d_{\tilde \W^c_f}(f^n(x),f^n(y))\le 2 d_{\tilde \W^c_{f_0}}(\phi^{-1} f^n(x),\phi^{-1} f^n(y))$$ and $$ d_{G}(\phi^{-1} f^n(x),\phi^{-1} f^n(y))\le d_{G}( f^n(x),f^n(y))+2\epsilon\le C_1\mu^nd(x,y)+2\epsilon\le 3\epsilon$$ for sufficiently large $n>N_\epsilon $, and $\epsilon$ can be taken to be arbitrarily small as $d_{C^1}(f,f_0)\to 0$. Since $d_{G}(\phi^{-1} f^n(x),\phi^{-1}f^n(y))\le 3\epsilon$ and $\phi^{-1} f^n(y)\in \tilde\W^c_{f_0}(\phi^{-1} f^n(x))$ the local geometry in $\tilde G$ gives a uniform bound $$C_3 d_{G}(\phi^{-1} f^n(x),\phi^{-1} f^n(y))\ge d_{\tilde \W^c_{f_0}}(\phi^{-1} f^n(x),\phi^{-1} f^n(y))$$ where $C_3>1$ depends only on $f_0$ and $G$.
 
 Therefore, we have $d_{\tilde \W^c_f}(f^n(x),f^n(y))<6C_3\epsilon $ for $n>N_\epsilon $. Pick a sufficiently small $\epsilon$, this implies that $d_{\tilde \W^c_f}(f^n(x),f^n(y)) $ is comparable to $d_X(f^n(x),f^n(y))$ for any $n>N_\epsilon $. On the other hand, we have $$d(f^n(x),f^n(y))\le C_1 \mu^n d(x,y)$$ and $$d_{\tilde \W^c_f}(f^n(x),f^n(y))\ge C_2 \nu^n d_{\tilde \W^c_f}(x,y)$$ for any $n\in \N$, since $0<\mu<\nu<1$, letting $n\to \infty$ gives a contradiction unless $x=y$. 
 
 Secondly, $u|_L$ is locally a homeomorphism, since local pieces of $\W^s_f$ are homeomorphic to pieces of $\W^s_{f_0}$ via projection in the $\W^c_f$ direction. Moreover, $u|_L$ is proper. Let $K$ be any compact set in $G^s$, then $u^{-1}(K)\cap L$ is a bounded set in $\tilde G$ since $\phi^{-1}$ is uniformly $C^0$-close to identity and $E^s_f$ is $C^0$-close to $E^s_{f_0}$. Therefore, $u(L)$ is both open and closed in $G^s$, so $u$ is also surjective.

 This shows that for any $x,y\in G$, $\tilde \W^s_f(x)\cap \tilde \W^c_f(y)=\{z\in \tilde \W^s_f(x): u(z)=u(y)\}$ is a unique point in $G$. Therefore, stable holonomies of $\tilde f$ are globally well-defined in $G$, and similarly, unstable holonomies of $\tilde f$ are globally well-defined in $G$. Passing to the quotient $G/\Gamma$, for any $su$-path $\gamma$, a lifting $\tilde \gamma$ can be globally extended to continuous family of paths from $\tilde \W^c(\tilde \gamma(0))$ to $\tilde \W^c(\tilde \gamma(1))$, since $\W^c_f(x_0)$ is simply connected, the projection $ G\to G/\Gamma$ restricted to $\tilde \W^c_f(x_0)$ is a homeomorphism, this continuous family of paths project to a continuous family of paths from $\W^c_f(x_0)$ to $\W^c_f(\gamma(1))$. This shows that $f$ has global holonomy on $\W^c_f(x_0)$.
 
\end{proof}

\end{document}
